\documentclass[10pt,a4paper]{amsart}
\usepackage[margin=19mm]{geometry}
\usepackage{amsmath,amssymb,mathtools}
\usepackage[scr=rsfs,cal=euler]{mathalfa}
\usepackage{xfrac}
\usepackage[unicode,hidelinks,bookmarks=false]{hyperref}
\newcommand{\ie}{i.e.\ }
\newcommand{\cf}{cf.\ }
\newcommand{\resp}{respectively\ }
\newcommand{\Subsection}{\subsection}
\providecommand{\nobreakdash}{\nobreak\hbox{-}}
\setlength{\emergencystretch}{2em}
\multlinegap0pt
\usepackage{mathtools}
\usepackage{amssymb}
\usepackage{xcolor}
\usepackage{dsfont}
\usepackage{float}
\usepackage[shortlabels]{enumitem}
\newtheorem{lemma}{Lemma}
\newtheorem{theorem}{Theorem}
\title{Chiral Polyhedra from $\mathrm{AGL}(1,q)$}
\author{Evan Angelone, Egon Schulte}
\date{Source: arXiv:2603.02543v1 (CC BY 4.0)}
\begin{document}
\maketitle

\noindent\emph{Source and licence.} Chiral Polyhedra from $\mathrm{AGL}(1,q)$. Version arXiv:2603.02543v1 (CC BY 4.0), distributed by arXiv under the Creative Commons Attribution 4.0 International licence, http://creativecommons.org/licenses/by/4.0/, as recorded in the arXiv licence metadata for this exact version and shown on the abstract page https://arxiv.org/abs/2603.02543v1. Full licence notice: \texttt{licenses/arxiv-2603.02543-CC-BY-4.0.txt}.\par\smallskip

\noindent\textit{Editorial scope.} Counted unit: the article's lemma conjlemma (source0.tex lines 394--413). The article's complete original proof of it is delivered with it (source0.tex lines 415--464, one \texttt{proof} environment of 50 lines). No mathematics is written, repaired or re-derived: the statement and the proof are the article's own lines, in the article's own notation, and no retained line is substituted or re-flowed.  Retained uncounted as the source-local context the proof needs: lines 182--185; lines 281--284; lines 291--294; lines 319--332.  Nothing was omitted from any retained span, so the package carries no omission record.  No retained line contains a citation, so the wrapper carries no bibliography and no bibliography entry.  No retained line contains a figure environment, an image-inclusion command or drawing code, so no image is delivered and the package declares no asset dependency.  Editorial formatting only, in addition: an AMS \texttt{amsart} preamble that restates the article's own declarations and macros used by the retained lines, namely the environments \texttt{lemma}, \texttt{theorem} on separate counters, together with the standard packages those lines need.  The retained environments print in this document as Lemma 1, Theorem 1 in order of appearance, whereas the article prints the same items with the numbering of its own full document.  The complete native source of the article as distributed in the arXiv e-print for this exact version is bundled unchanged as \texttt{sources/195627/source0.tex}.
\begin{lemma}
\label{invgrauts}
Let $\rho:=\gamma(a,b,\sigma)$ be a group automorphism of $\mathrm{AGL}(1,q)$, where $a\in\mathbb{F}_q^*$, $b\in\mathbb{F}_q$, and $\sigma\in\rm{Aut}(\mathbb{F}_q)$ is not the identity. If $\rho$ is an involution, then $\sigma$ is also an involution (thus $l$ is even and $\sigma=\varphi^{l/2}$)
and $aa^{\sigma} = 1$ and $ab^\sigma +b=0$. \end{lemma}

\begin{equation}
\label{sig2}
\sigma_{2}\coloneqq \alpha(a,b),\;\tau\coloneqq \alpha(-1,c),
\end{equation}

\begin{equation}
\label{sig1}
\sigma_{1}\coloneqq \tau\sigma_2^{-1} = \alpha(-1,c)\,\alpha(a^{-1},-a^{-1}b)=\alpha(-a^{-1},a^{-1}b+c) .
\end{equation}

\begin{theorem}
\label{gamprops}
Let $q$ be an odd prime power, $q\neq 3$, and let $\Gamma\coloneqq \langle \sigma_1,\sigma_2\rangle$ where the generators $\sigma_1$ and $\sigma_2$ are defined as in \eqref{sig2} and~\eqref{sig1}, with $a\in\mathbb{F}_q^*$, $b,c\in\mathbb{F}_q$, $a\neq \pm 1$, and $c\neq 2b/(1-a)$. Let $\mathbb{F}_q^*=\langle\gamma\rangle$ and $a=\gamma^m$, $1\leq m<q-1$. Further, let $T_\Gamma\coloneqq T\cap \Gamma$ and $H_\Gamma\coloneqq H\cap \Gamma\, (\cong T_\Gamma\rtimes C_2)$. Then, \\[.03in]
(a)\ $\Gamma$ is either the group of a chiral polyhedron or the rotation subgroup of a regular polyhedron $\mathcal{P}$, in either case with distinguished generators $\sigma_1,\sigma_2$. \\[.03in]
(b)\ $\mathcal{P}$ is of type $\{s,t\}$, where $s=o(-a^{-1})$ and $t=o(a)$ are the orders of $-a^{-1}$ and $a$ in $\mathbb{F}_q^*$, respectively. \\[.03in]
(c)\ $\Gamma \cong T_\Gamma \rtimes \langle \sigma_2 \rangle$ if $\gcd(q-1,m)\,|\,r$ and $\Gamma \cong H_\Gamma \ltimes \langle \sigma_2 \rangle$ if $\gcd(q-1,m)\!\!\not|\,r$.\\[.03in]
(d)\ $H_\Gamma = N(\tau)$, where $N(\tau)$ denotes the normal closure of $\tau=\sigma_1\sigma_2$ in $\Gamma$.
\\[.03in]
(e)\ $T_\Gamma$ contains all translations by elements of the subgroup $M_\Gamma \coloneqq  \langle\, (a^k-1)g\mid k\geq 0\rangle$, 
with $g\coloneqq c-\tfrac{2b}{1-a}$, of the additive group $\mathbb{F}_q$ whose rank is at least 
$\lceil \log_{p} ((q-1)/\gcd(q-1,m))\rceil$. \\[.03in]
(f)\ $T_{\Gamma}\cong \mathbb{F}_q$, at least if $m$ and $q-1$ are relatively prime or $q$ itself is a prime.\\[.03in]
(g)\ $\Gamma=\mathrm{AGL}(1,q)$ if and only if either $\gcd(q-1,m)=1$ (that is, $\mathbb{F}_q^*=\langle a\rangle$) or $q\equiv 3\bmod 4$ and ${\gcd(q-1,m)=2}$.
\end{theorem}

\begin{lemma}
\label{conjlemma}
Let $\rho:=\gamma(g,h,\sigma)$ be an involutory group automorphism of $\mathrm{AGL}(1,q)$, where $\sigma$ is an involution in $\rm{Aut}(\mathbb{F}_q)$ and $g\in\mathbb{F}_q^*$, $h\in\mathbb{F}_q$ with $g\overline{g}= 1$, $g\overline{h}+h=0$ (see Lemma~\ref{invgrauts}). Let $\Gamma=\langle\sigma_2,\tau\rangle$ be as in Theorem~\ref{gamprops}, with $\sigma_{2}=\alpha(a,b)$ and $\tau=\alpha(-1,c)$ as in (\ref{sig2}). Suppose $\rho$ acts by conjugation as a group automorphism on $\Gamma$, that is, $\rho\,\Gamma\rho=\Gamma$. \\[.03in]
(a)\ If $\rho \sigma_2\rho = \sigma_2^{-1}$, then $\overline{a}=a^{-1}$ and thus $a^{p^{l/2} +1}=1$.\\[.03in]
(b)\ If $\rho \sigma_2\rho = \sigma_2^{-1}$ and $\rho\tau\rho=\tau$, then $\overline{a}=a^{-1}$ and $g,h$ are the unique solutions of the nonsingular linear system
\begin{equation}
\label{linsystem}
\left\{\begin{array}{rcrcr}
\overline{b} g &+& (1-\overline{a})h &=& -\overline{a}b \\
\overline{c} g &+& 2h &= &c,
\end{array}\right.
\end{equation}
where for simplicity the variables of the system were also denoted by $g$ and $h$. The solutions to \eqref{linsystem} are given by 
\begin{equation}
\label{systemsols}
g=-\overline{a}\,(\overline{D}/{D}),
\;\;h=(\overline{b}c+\overline{a}b\overline{c})/D,
\end{equation}
where $D:=2\overline{b}-\overline{c}(1-\overline{a})$ is the (non-vanishing) determinant of the linear system.
\end{lemma}

\begin{proof}
The conjugate of $\sigma_2$ by $\rho$ is given by
\begin{equation}
\label{conjrho2}
\begin{array}{rcl}
\rho\sigma_2\rho(x)\,=\,\rho\sigma_2 (g\overline{x}+h)\!\!&=&\!\!\rho(a(g\overline{x}+h)+b)\\[.03in]
\!\!&=&\!\!g\;\overline{(ag\overline{x}+ah+b)}+h \\[.03in]
&=&\!\! (g\overline{g})\overline{a}x+g\overline{a}\overline{h}+g\overline{b} +h
\,=\,\overline{a}x+g\overline{a}\overline{h}+g\overline{b} +h.
\end{array} 
\end{equation}
If this conjugate is to coincide with $\sigma_2^{-1}=\alpha(a^{-1},-a^{-1}b)$, then 
\[a^{-1}x-a^{-1}b = \overline{a}x+g\overline{a}\overline{h}+g\overline{b} +h\]
for all $x\in\mathbb{F}_q$ and therefore
\begin{equation}
\label{consigma}
\overline{a}=a^{-1},\;\, g\overline{a}\overline{h}+g\overline{b} +h=-a^{-1}b.
\end{equation}
In particular this proves the first part of the lemma. Note that $\overline{a}=a^{-1}$ just means $a^{p^{l/2} +1}=1$. Further, by assumption, $g\overline{h}+h=0$, so using $\overline{a}=a^{-1}$ allows us to rewrite the second equation in \eqref{consigma} as 
\begin{equation}
\label{chireq1}
\overline{b}g +(1-\overline{a})h=-\overline{a}b,
\end{equation}
which is a the first linear equation of \eqref{linsystem} in $g$ and $h$.

Similarly, the conjugate of $\tau$ by $\rho$ is given by
\begin{equation}
\label{conjtau}
\begin{array}{rcl}
\rho\tau\rho(x)\,=\,\rho\tau(g\overline{x}+h)\!\!&=&\!\!\rho(-(g\overline{x}+h)+c)\\[.03in]
\!\!&=&\!\! g\,\overline{(-g\overline{x}-h+c)}+h \\[.03in]
&=&\!\!-(g\overline{g})x-\!g\overline{h}+g\overline{c} +h
\,=\,-x-g\overline{h}+g\overline{c} +h.
\end{array} 
\end{equation}
If this conjugate is to coincide with $\tau$ itself, then $c=-g\overline{h}+g\overline{c} +h$ and therefore 
\begin{equation}
\label{chireq2}
\overline{c}g + 2h = c,
\end{equation}
which is the second linear equation of \eqref{linsystem} in $g$ and $h$. Here, we once again used that $g\overline{h}+h=0$. 

Therefore, if both $\rho \sigma_2\rho = \sigma_2^{-1}$ and $\rho\tau\rho=\tau$ hold, then the equations in (\ref{chireq1}) and (\ref{chireq2}) give the linear system of (\ref{systemsols}) in $g$ and $h$ with determinant $D=2\overline{b}-\overline{c}(1-\overline{a}) = \overline{2b-c(1-a)}$. 

Now recall from (\ref{sig2}) that our initial conditions on the parameters $a,b,c$ for the generators $\sigma_2,\tau$ of $\Gamma$ included the condition $c\neq 2b/(1-a)$ which ensured the intersection property for $\Gamma$. This shows that the linear system is nonsingular and has a unique solution in $g$ and $h$. Bearing in mind that $\overline{a}=a^{-1}$, the solutions can be expressed as 
\[ g=\frac{-2\overline{a}b-c(1-\overline{a})}{2\overline{b}-\overline{c}(1-\overline{a})}=
-\overline{a}\,\frac{\;\overline{D}\;}{\,D\,},
\;\;\;\; h=\frac{\overline{b}c+\overline{a}b\overline{c}}{2\overline{b}-\overline{c}(1-\overline{a})}=\frac{\overline{b}c+\overline{a}b\overline{c}}{D}.\]
Then we can check that, indeed, $g\overline{g}=1$ and $g\overline{h}+h=0$, as required by our assumptions. 
\end{proof}

\end{document}
